\section{Dissipation-clock continuation without an endpoint theorem}
\label{dc:section}

\paragraph{Scope and review status.}
This section gives full component derivations for the original unforced
three-dimensional equation, not an arbitrary-data regularity proof.  The
strict-absorption consequence reconstructs and sharpens the separate
\texttt{dissipation\_budget\_continuation.tex} supplement: the useful Serrin
pair is $(3,9)$, directly supplied by cubic dissipation, rather than $(4,6)$.
The sharper estimate and the nonlinear barriers below are author-checked
additions awaiting an independent mathematical audit; they are not promoted
claims in the dependency graph.  The nonendpoint testing mechanism is
classical.  No novelty claim is made.

\paragraph{Inputs and conventions.}
Fix $\nu>0$, a real divergence-free Schwartz datum $u_0$ on $\mathbb R^3$,
and its maximal classical branch $u$ on $[0,T_*)$.  The imported local
package, assembled earlier from Tao's local and maximal theory
\cite{Tao2013}, gives $u\in C^j([0,T];H^k)$ for every $T<T_*$ and all
$j,k$, and $\|u(t)\|_{H^1}\to\infty$ at a finite $T_*$.  We use the energy
identity and the integrated cubic pressure balance already proved for this
branch.  Put
\begin{align*}
 E(t)&=\|u(t)\|_2^2,&Y(t)&=\|\nabla u(t)\|_2^2,&Z(t)&=\|\Delta u(t)\|_2^2,\\
 X(t)&=\|u(t)\|_3^3,&
 D_3(t)&=\int_{\mathbb R^3}|u|
          \bigl(|\nabla u|^2+|\nabla |u||^2\bigr)\,dx,&
 B(t)&=\int_0^tD_3(s)\,ds.
\end{align*}
Here $\nabla|u|$ is the weak gradient, taken as zero almost everywhere on
$\{u=0\}$; equivalently the second integrand is
$|(\nabla u)^Tu|^2/|u|$, with zero value on that set.  With
$P_3=\int p\,u\cdot\nabla|u|$, the two input balances are
\begin{equation}\label{dc:balances}
 E(t)+2\nu\int_0^tY=E(0),\qquad
 \frac{X(t)}3+\nu B(t)=\frac{X(0)}3+\int_0^tP_3.
\end{equation}
Let $S$ be any constant in $\|f\|_6\le S\|\nabla f\|_2$ on
$\dot H^1(\mathbb R^3)$.  The same constant applies to finite-dimensional
vector and tensor fields by the weak norm-gradient inequality.  All norms
below are spatial.  No assumption about the value of a continuation norm is
part of these inputs.

\begin{lemma}[The cubic dissipation supplies a nonendpoint Serrin norm]
\label{dc:l9}
On every compact classical interval,
\begin{equation}\label{dc:l9-bound}
 \|u\|_9^3\le \frac{9S^2}{8}D_3.
\end{equation}
\end{lemma}
\begin{proof}
Set $r=|u|$ and $f=r^{3/2}$.  The local package gives $u\in L^3\cap
L^\infty$ and $\nabla u\in L^2$.  Thus $f\in H^1$ and
$\nabla f=(3/2)r^{1/2}\nabla r$ almost everywhere, including the zero set.
This follows either from the Sobolev chain rule with bounded truncations or
by smoothing the norm and passing to the limit.  Since
$|\nabla r|\le |\nabla u|$ almost everywhere,
\[
 \|u\|_9^3=\|f\|_6^2
 \le S^2\|\nabla f\|_2^2
 =\frac{9S^2}{4}\int r|\nabla r|^2
 \le\frac{9S^2}{8}D_3.
\]
There is no division by $r$ in this argument.
\end{proof}

\begin{proposition}[Enstrophy in the dissipation clock]
\label{dc:clock}
For every $t<T_*$,
\begin{equation}\label{dc:clock-differential}
 Y'(t)+\nu Z(t)\le \frac{4S^3}{3\nu^2}D_3(t)Y(t),
\end{equation}
and consequently
\begin{equation}\label{dc:clock-bound}
 Y(t)+\nu\int_0^t Z(s)\,ds
 \le Y(0)\exp\!\left(\frac{4S^3}{3\nu^2}B(t)\right).
\end{equation}
In particular, for any finite $H>0$, a uniform finite bound on
$B(t)$ for $t<\min(H,T_*)$ implies $T_*>H$.  Conversely, $T_*>H$ implies
$B(H)<\infty$.
\end{proposition}
\begin{proof}
Test the equation against $-\Delta u$.  The pressure cancels because
$\operatorname{div}u=0$.  Spatial integration by parts is justified by the
local Sobolev package, for example by cutoffs followed by their removal.
H\"older, interpolation and Sobolev give
\begin{align*}
 \frac12Y'+\nu Z
 &\le \|u\|_9\|\nabla u\|_{18/7}\|\Delta u\|_2\\
 &\le S^{1/3}\|u\|_9Y^{1/3}Z^{2/3}.
\end{align*}
Indeed $1/9+7/18+1/2=1$ and
$\|\nabla u\|_{18/7}\le\|\nabla u\|_2^{2/3}
\|\nabla u\|_6^{1/3}$; moreover
$\|\nabla u\|_6\le S\|D^2u\|_2=S\|\Delta u\|_2$, the last equality
being the Fourier Hessian identity.  For $a,z\ge0$,
\[
 2az^{2/3}\le\nu z+\frac{32}{27\nu^2}a^3.
\]
For $a>0$ this is the exact maximum of $2az^{2/3}-\nu z$, attained at
$z=(4a/(3\nu))^3$; at $a=0$ it is immediate.  Therefore
\[
 Y'+\nu Z\le\frac{32S}{27\nu^2}\|u\|_9^3Y
 \le\frac{4S^3}{3\nu^2}D_3Y,
\]
where $(32/27)(9/8)=4/3$.  Put $c=4S^3/(3\nu^2)$.  Multiplication by
$e^{-cB(t)}$ and integration yield
\[
 e^{-cB(t)}Y(t)+\nu\int_0^t e^{-cB(s)}Z(s)\,ds\le Y(0).
\]
Since $B$ is nondecreasing, multiplication by $e^{cB(t)}$ implies
\eqref{dc:clock-bound}.  This also handles $Y(0)=0$ without division.

If $T_*\le H<\infty$ and $B$ is bounded below $T_*$, then
\eqref{dc:clock-bound} bounds $Y$ there.  The energy identity bounds $E$,
contradicting the local $H^1$ blow-up alternative.  This includes the case
$T_*=H$.  Conversely the local package makes $D_3$ integrable on $[0,H]$
whenever $T_*>H$.  No endpoint $L^\infty_tL^3_x$ theorem, backward
uniqueness, comparison flow, or quotient-minimizer regularity is used.
\end{proof}

\begin{corollary}[Strict signed pressure absorption: explicit direct suffix]
\label{dc:strict}
Fix $H>0$.  Suppose there are $0\le\theta<1$ and a finite $A\ge0$, chosen
before $t$, such that for every $t<\min(H,T_*)$,
\begin{equation}\label{dc:strict-hyp}
 \int_0^tP_3\le\theta\nu B(t)+A.
\end{equation}
Then $T_*>H$, and on that interval, with $R=X(0)+3A$,
\begin{equation}\label{dc:strict-bound}
 Y(t)+\nu\int_0^tZ
 \le Y(0)\exp\!\left(\frac{4S^3R}{9(1-\theta)\nu^3}\right).
\end{equation}
\end{corollary}
\begin{proof}
Subtract \eqref{dc:strict-hyp} from the cubic balance, retaining both
nonnegative terms:
\[
 X(t)+3(1-\theta)\nu B(t)\le R.
\]
Thus $B(t)\le R/[3(1-\theta)\nu]$.  Apply
Proposition~\ref{dc:clock}.  In particular the original fixed-cutoff
HIGH-PRESSURE hypothesis, together with the proved low-pressure bound,
implies this corollary with $A=A_{\rm high}+L_J(H)$.  The cutoff and
remainder quantifiers are unchanged.  Strictness is used only to bound
$B$, not to invoke an endpoint theorem.
\end{proof}

\begin{theorem}[A single finite barrier in accumulated dissipation]
\label{dc:barrier}
Fix $H>0$ and an integer cutoff $J$.  Write $Q_J$ for high-output pressure
work and suppose the established low-output estimate is
\[
 \left|\int_0^t(P_3-Q_J)\right|\le L_J(H)<\infty
 \quad\text{for every }t<\min(H,T_*).
\]
Suppose there are $b>0$ and $\varepsilon>0$, chosen from the inputs before
$t$, with the following property: if $\tau_b<\min(H,T_*)$ is the first
time at which $B(\tau_b)=b$, then
\begin{equation}\label{dc:crossing-hyp}
 \int_0^{\tau_b}Q_J
 \le\nu b-\frac{X(0)}3-L_J(H)-\varepsilon.
\end{equation}
Then $B(t)<b$ below $\min(H,T_*)$, $T_*>H$, and the right-hand side of
\eqref{dc:clock-bound} is at most $Y(0)e^{4S^3b/(3\nu^2)}$.
\end{theorem}
\begin{proof}
$B$ is continuous, nondecreasing and zero at time zero.  If it reaches
$b$ before $\min(H,T_*)$, continuity gives a first such time.  At that
time \eqref{dc:balances}, the low-output bound and
\eqref{dc:crossing-hyp} imply $X(\tau_b)/3\le-\varepsilon$, impossible.
Consequently $B<b$ throughout the open interval.  Apply
Proposition~\ref{dc:clock}; no assertion at the unconstructed singular
endpoint is required.
\end{proof}

\begin{corollary}[Nonlinear deficits and vanishing absorption margins]
\label{dc:nonlinear}
In Theorem~\ref{dc:barrier}, suppose instead that a finite $A_{\rm high}\ge0$
and a function $\Phi:[0,\infty)\to\mathbb R$, both fixed from the inputs,
satisfy
\begin{equation}\label{dc:nonlinear-hyp}
 \int_0^tQ_J\le\nu B(t)-\Phi(B(t))+A_{\rm high}
 \quad(t<\min(H,T_*)).
\end{equation}
It suffices that there be one finite $b>0$ for which
\begin{equation}\label{dc:barrier-value}
 \Phi(b)>K,\qquad K=\frac{X(0)}3+A_{\rm high}+L_J(H).
\end{equation}
No monotonicity, continuity, or coercivity of $\Phi$ is required for this
single-barrier conclusion.  In particular, with $b_0,\eta>0$, the choice
\begin{equation}\label{dc:log-deficit}
 \Phi(B)=\eta\nu b_0\log(1+B/b_0)
\end{equation}
gives continuation and the explicit bound
\begin{equation}\label{dc:log-bound}
 B(t)\le b_0\left[\exp\!\left(\frac{K}{\eta\nu b_0}\right)-1\right].
\end{equation}
\end{corollary}
\begin{proof}
At the hypothetical first crossing use
$\varepsilon=\Phi(b)-K>0$ in Theorem~\ref{dc:barrier}.
For the logarithmic choice, either this theorem applies at any $b$ larger
than the right-hand side of \eqref{dc:log-bound}, or one can directly
subtract \eqref{dc:nonlinear-hyp} from the cubic balance to obtain
$X(t)/3+\Phi(B(t))\le K$ and invert the increasing logarithm.
The positive constants ensure a finite barrier even when $K=0$.
\end{proof}

\paragraph{What the logarithm does and does not remove.}
Here $\Phi(B)/B\to0$ as $B\to\infty$, so a fixed positive fractional
margin below viscosity is not necessary in this sufficient-condition
framework.  An example of a stronger, pointwise hypothesis implying
\eqref{dc:nonlinear-hyp} is
\[
 Q_J(t)\le\left(\nu-\frac{\eta\nu}{1+B(t)/b_0}\right)D_3(t)+a(t),
 \qquad \int_0^t a(s)\,ds\le A_{\rm high},
\]
by the ordinary chain rule for the absolutely continuous $B$.  Neither
this inequality nor \eqref{dc:nonlinear-hyp} has been proved for arbitrary
Navier--Stokes data.  Choosing $\Phi$ is not a way to establish the
pressure estimate that contains it.

\begin{remark}[Exact scalar boundary; not a PDE counterexample]
\label{dc:scalar}
A bounded deficit whose supremum stays below the available budget does not
control $B$ through the cubic balance alone.  For a nonnegative bounded
$C^1$ function $\Phi$ with $\Phi(0)=0$, choose $K>\sup\Phi$, and for
$0\le t<1$ set
\[
 B(t)=-\log(1-t),\quad D(t)=\frac1{1-t},\quad
 X(t)=3\bigl(K-\Phi(B(t))\bigr),\quad
 P(t)=\bigl(\nu-\Phi'(B(t))\bigr)D(t).
\]
Then $X>0$, $X'/3+\nu D=P$ and
$\int_0^tP=\nu B(t)-\Phi(B(t))$, but $B(t)\to\infty$.
These are scalar functions, not a velocity field, pressure, or solution of
Navier--Stokes.  This only rejects that scalar inference.  In particular
$\theta=1$ with a constant remainder is not justified by the direct
suffix.  The separately proved bare $L^\infty_tL^3_x$ criterion still
retains its endpoint theorem.
\end{remark}

\paragraph{Remaining quantified mathematical obligation.}
For every viscosity, every admissible datum and every finite horizon, one
must still \emph{derive}, from the original equation, the fixed-cutoff
HIGH-PRESSURE estimate, a barrier such as \eqref{dc:crossing-hyp}, or some
other sufficient budget.  None is supplied by the energy identity or by
merely introducing the dissipation clock.  At the existential level the
barrier formulation is also equivalent to continuation: on a branch already
known to extend past $H$, choose $b>B(H)$ and its first-crossing condition
is vacuous.  That converse is a scope observation, not an a priori
construction.  No terminal or open estimate is discharged by this section.
